\section{Сохранённый расчёт на прежних условиях}
\label{app:calculation-history}

База расчёта этого приложения использует сервер llama.cpp с Q4, прежнюю тарификацию API и условия хранения run 16.
Числа сохраняют результат той модели; текущие границы применимости с сервером vLLM и измеренной долей кэша приведены в разделе~\ref{sec:economics}.
Вариант RTX 5090 как сервер в текущую модель не входит.

\subsection{Сохранённая база расчёта с llama.cpp Q4}
\label{sec:economics-legacy}

Следующие таблицы и рисунки сохраняют расчёт на серверной кривой llama.cpp Q4 и прежних предположениях о кэше API.
Эта база не задаёт текущие границы применимости vLLM FP8; границы в начале раздела ожидают пересчёта.
Прежние значения сравнений и цен сохранены как именованная чувствительность.

Все числа этого раздела посчитаны движком расчётов из репозитория статьи. Каждое входное значение там помечено как
замер, как цена с источником или как оценка. Профиль нагрузки тяжёлый: 2\,500 вызовов в день на человека, 21 рабочий
день в месяц, восьмичасовой рабочий день. На открытую модель уходит доля $s$ этих вызовов (раздел~\ref{sec:workload}).
Модель Qwen3.8-27B, сервер считает в тех же 4-битных весах, что и клиент. Цены API взяты у поставщиков этой же модели: на нашем профиле
1000 вызовов у медианного поставщика стоят \$9.68, у самого дешёвого \$1.34. Генерация на вызов равна ответу по
журналам, 512 токенов, в которые рассуждения модели уже входят (раздел~\ref{sec:workload}), а прошлый ответ модели
оплачивается как новый ввод. Если прибавить рассуждения к ответу ещё раз, получается верхняя граница генерации:
ёмкость M5 Max падает с 838 до 605 вызовов в день, API дорожает на 5--20\%, соотношение с медианным поставщиком
меняется мало (при доле 50\% гибрид \$57 против \$268), а разрыв с самым дешёвым сокращается: \$57 против \$42 при
доле 50\% и в 1.9 раза при доле 20\%.

\textbf{Годовой счёт.} Для масштаба мы оценили по ценам API октября 2026 года сохранившиеся журналы
собственных сессий одного пользователя: 1.53 миллиона вызовов агентов, у Codex с 2025-11-18, у Claude с 2026-05-07,
у Muse с 2026-08-06 (таблица~\ref{tab:owner}). По спискам цен вызовы стоили бы \$113--124 тысячи. Вызов Claude обходится
в среднем в 16 центов, вызов Codex в 6--7 центов; вызов в гибриде с Qwen3.8-27B при доле 20\% стоит 0.64--0.82 цента. По индексу
Artificial Analysis на моделях уровня Qwen3.8-27B и ниже прошло 1.4\% вызовов, ещё 15.6\% на моделях, которые не
выше Qwen3.8-27B на стандартном усилии рассуждения. Вместе это не больше 17\% вызовов (с каким усилием они шли, журналы не
фиксируют), близко к доле $s = 20\%$ наших расчётов, но только 5\% счёта: такие вызовы шли на дешёвых моделях. Основные деньги приходятся на передовые модели.
Классы здесь определены по модели; какие задачи передовых моделей можно отдать открытой модели, по журналам сказать
нельзя. Ещё 6.2 тысячи вызовов в журналах уже выполнены на открытых моделях локально на Маке.

\begin{table}[t]
\centering
\caption{Сохранившиеся журналы одного пользователя по ценам API октября 2026 года, по классам моделей.
Срез по моделям содержит 1\,530\,739 вызовов; отличие от выборки нагрузки разобрано в разделе~\ref{sec:workload}. Класс задан
индексом Artificial Analysis на той же версии индекса, что и у Qwen3.8-27B. Диапазоны отражают изменения цен за год
и неизвестные в журналах величины, например смысл кэша у Muse. Запись в кэш у Codex журналы показывают нулевой на
каждом вызове, и она так и посчитана. Данные и расчёт: \texttt{data/owner-model-usage-2026-10-30.json},
\texttt{data/year\_bill.py}.}
\label{tab:owner}
\footnotesize
\begin{tabular}{lrrr}
\toprule
Класс моделей & Вызовов & Доля & Счёт \\
\midrule
Не выше Qwen3.8-27B & 21\,935 & 1.4\% & \$0.6 тыс. \\
Не выше на стандартном усилии рассуждения & 238\,266 & 15.6\% & \$4.8--5.1 тыс. \\
Выше Qwen3.8-27B & 1\,248\,111 & 81.5\% & \$107--118 тыс. \\
Без индекса & 13\,157 & 0.9\% & \$0.3 тыс. \\
Открытые модели локально на Маке & 6\,178 & 0.4\% & --- \\
Без цены в списке & 3\,092 & 0.2\% & --- \\
\midrule
Всего & 1\,530\,739 & 100\% & \$113--124 тыс. \\
\bottomrule
\end{tabular}
\end{table}

\textbf{Мак уже есть.} Пусть у пользователя уже есть машина класса Max. Тогда гибрид добавляет к его расходам
серверный prefill, хранение копий сессий на сервере, трафик, электричество клиента и обслуживание сервиса. Серверный
prefill мы оплачиваем по часовой цене B300, \$9.50 в час при загрузке 70\%, трафик по \$5 за терабайт, обслуживание
оцениваем в \$5 на человека в месяц. Хранение считаем по политике, которую допускает протокол: сервер дописывает в
копию сессии только новые позиции, держит копии в оперативной памяти в рабочие часы и на NVMe в остальное время, а токены вызовов, прочитанных клиентом, дочитывает при следующем оффлоаде. Гигабайт в месяц мы оцениваем в \$0.50 в оперативной памяти и \$0.10 на NVMe. Для M5 Max при доле 20\% хранение стоит \$5.6 в месяц, меньше, чем сам prefill (\$54.9): хранить дешевле.\footnote{Снимок run 16: \texttt{person\_cost\_s20.csv}, строка M5 Max 64 GB, Mac already owned, команда 50; поля \texttt{hosting\_month} и \texttt{compute\_month}.} Осторожная политика, при которой копии весь месяц лежат в оперативной памяти, а сессия на каждом
вызове перекачивается на GPU, обходится в \$16.5. Таблица~\ref{tab:econ} показывает итог для разных клиентов.

При доле 20\% пользователь отправляет на открытую модель 500 вызовов в день, и клиенты класса Max почти все их
обслуживают сами. Гибрид на клиентах классов Max и Ultra обходится в \$67--73 в месяц. Эта месячная цена в 4.8--5.2 раза дороже, чем у самого дешёвого поставщика API
(\$14), и в 1.4--1.5 раза дешевле, чем у медианного (\$102). При осторожной политике
хранения гибрид в 5.8--6.0 раза дороже самого дешёвого поставщика и в 1.2 раза дешевле медианного.

При доле 50\% на открытую модель уходит 1\,250 вызовов в день, и здесь проявляется предел клиента. Генерация
занимает время, и за восьмичасовой день M5 Max с четырьмя потоками успевает обслужить около 1010 вызовов. Остальные мы
отправляем самому дешёвому поставщику API и включаем их стоимость в цену гибрида. Выходит \$107--160 в месяц: на 204--354\% дороже самого
дешёвого API (\$35) и в 1.6--2.4 раза дешевле медианного (\$254). Это соотношение складывает два
эффекта: гибрид и переход к самому дешёвому поставщику для вызовов сверх ёмкости клиента. Если эти вызовы оставить у
медианного поставщика, гибрид на клиентах класса Max дешевле его только в 1.2--1.6 раза. Клиенты класса Ultra обслуживают все 1\,250 вызовов
сами и обходятся в \$150, на 327\% дороже самого дешёвого API при хранении по политике протокола и в \$176--177 при осторожной политике.

\begin{table}[t]
\centering
\caption{Гибрид на человека в месяц, тяжёлый профиль (2\,500 вызовов в день). ‡ M3: промежуточные многопоточные опоры PRI~4, повтор при PRI~31 ожидается. M5: приоритет не записан; запускающий процесс показал PRI~31 при последующей проверке. Вызовы сверх ёмкости клиента оплачены
по цене самого дешёвого поставщика API. Хранение копий сессий на сервере по политике протокола; осторожная политика
хранения добавляет \$10--29. Ёмкость рассчитана для четырёх потоков генерации за восьмичасовой день
по измеренным кривым M3 Max и M5 Max с 64~ГБ. M5 Max с 128~ГБ использует измеренные кривые той же SoC M5 Max с 64~ГБ;
M5 Ultra масштабирует эти кривые по вычислениям и полосе памяти. Перенос на M5 Max с 128~ГБ и M5 Ultra
не подтверждён отдельными измерениями. M4 Max оценён по вычислительной мощности и полосе памяти.}
\label{tab:econ}
\footnotesize
\begin{tabular}{lrrrr}
\toprule
Клиент & Ёмкость, & $s = 20\%$, & $s = 50\%$, & $s = 50\%$, \\
 & вызовов в день & Мак уже есть & Мак уже есть & Мак куплен \\
\midrule
M3 Max 128 ГБ & 634 & \$73.4 & \$107.2 & \$154.7 \\
M4 Max 128 ГБ & 1\,277 & \$71.3 & \$160.0 & \$207.4 \\
M5 Max 64 ГБ & 1\,006 & \$71.3 & \$138.4 & \$173.6 \\
M5 Max 128 ГБ & 979 & \$71.7 & \$136.5 & \$193.3 \\
M5 Ultra 96 ГБ & 1\,503 & \$67.5 & \$150.5 & \$225.0 \\
M5 Ultra 256 ГБ & 1\,582 & \$67.5 & \$150.5 & \$296.7 \\
\midrule
API, медианный поставщик & & \$101.7 & \$254.2 & \$254.2 \\
API, самый дешёвый поставщик & & \$14.1 & \$35.2 & \$35.2 \\
\bottomrule
\end{tabular}
\end{table}

\textbf{Доля открытой модели.} Чем больше доля, тем больше гибрид экономит против медианного поставщика и тем раньше
упирается в ёмкость клиента. Пользователю с 800 вызовами в день ёмкости хватает и при доле 50\%. Долю 100\% мы
приводим как верхнюю границу: тяжёлому пользователю, который все вызовы отправляет открытой модели, понадобилось бы примерно три M5 Max с 64~ГБ, около четырёх M3 Max или около полутора Ultra.

\textbf{Покупка клиента.} Если машину покупают ради гибрида, к цене добавляется доплата за неё сверх обычного
ноутбука. Для M5 Max с 64~ГБ это \$2\,600, которые за четыре года при остаточной стоимости 35\% дают \$35 в месяц.
Тогда гибрид при доле 50\% стоит \$174 против \$35 у самого дешёвого API и \$254 у медианного. При доле 20\% ежемесячная цена купленного клиента выше медианного
поставщика: \$107 против \$102. Рисунок~\ref{fig:payback} показывает, за сколько
месяцев экономия против медианного поставщика возвращает доплату и машину целиком, когда вызовы сверх ёмкости
клиента уходят к самому дешёвому поставщику. Тогда при доле 50\% доплата возвращается за 22 месяца, Мак целиком за 44; при доле 20\% за 86 и 168 месяцев. Если вызовы сверх ёмкости оставить у медианного поставщика, при доле 50\%
доплата возвращается за 36 месяцев, Мак целиком за 72. Против самого дешёвого поставщика гибрид дороже,
и доплата не окупается. Доплату обычно платят к рабочему ноутбуку, который
специалисту нужен и так, и кроме экономии на вызовах она покупает память под модель и контекст, которые иначе не
помещаются, полосу памяти для быстрой генерации, несколько агентов одновременно и работу без поставщика. Против
самого дешёвого API доплату может окупить только эта польза: при доле 50\% гибрид с купленным M5 Max дороже его на \$138 в месяц. Сколько полезной работы даёт эта разница, покажет эксплуатация.

\begin{figure}[t]
\centering
\begin{tikzpicture}
\begin{axis}[
  ybar, bar width=11pt, width=0.62\textwidth, height=50mm, ymin=0, ymax=180,
  symbolic x coords={20,50,100}, xtick=data, xticklabels={{$s = 20\%$},{$s = 50\%$},{$s = 100\%$}},
  ylabel={месяцев}, enlarge x limits=0.25,
  nodes near coords={\pgfmathprintnumber[fixed,precision=0]{\pgfplotspointmeta}},
  nodes near coords style={font=\scriptsize},
  tick label style={font=\footnotesize}, label style={font=\footnotesize},
  legend style={font=\footnotesize,draw=none,at={(0.5,-0.2)},anchor=north,legend columns=2,column sep=6pt}]
\addplot[fill=white,draw=black] table[x=share,y=premium_months,col sep=comma] {data/payback.csv};
\addlegendentry{доплата за M5 Max 64~ГБ, \$2\,600}
\addplot[pattern=north east lines,draw=black] table[x=share,y=whole_months,col sep=comma] {data/payback.csv};
\addlegendentry{машина целиком, \$5\,099}
\draw[dashed] ({rel axis cs:0,0}|-{axis cs:20,48}) -- node[pos=0.5,above,font=\scriptsize]{четыре года} ({rel axis cs:1,0}|-{axis cs:20,48});
\end{axis}
\end{tikzpicture}
\caption{Окупаемость клиента против медианного поставщика API: доплата за M5 Max 64~ГБ сверх обычного ноутбука и
машина целиком, тяжёлый профиль. Вызовы сверх ёмкости клиента уходят к самому дешёвому поставщику; если оставить их
у медианного, при доле 50\% доплата возвращается за 36 месяцев, машина целиком за 72. Месяцы без учёта остаточной
стоимости. Против самого дешёвого поставщика гибрид дороже на \$57--103 в месяц, и покупка не окупается.}
\label{fig:payback}
\end{figure}

\textbf{Команды от 5 до 5000 человек.} Рисунок~\ref{fig:teams} сравнивает гибрид со своим серверным кластером при одинаковых требованиях во всём диапазоне команд.
Сравнение с API при долях 20 и 50\% сохраняет отдельный фиксированный аппаратный пример B300. Узел
B300 покупается целиком в этом аппаратном примере, поэтому для малой команды свой кластер очень дорог: на пять человек при лёгкой нагрузке с одним узлом \$3\,373 на человека в
месяц с учётом хранения. С ростом команды цена падает: на 500 человек фиксированный кластер B300 дороже медианного поставщика API при обеих долях
(\$674 против \$102 при доле 20\%, \$1685 против \$254 при доле 50\%), на 50 человек кластер B300 тоже дороже. Самого дешёвого
поставщика API свой кластер не догоняет ни при каком размере команды до 5000 человек. Гибрид от размера команды не зависит, потому что серверный prefill здесь покупается долями у общего поставщика по
часовой цене B300 при загрузке 70\%; команде со своим выделенным сервером такая цена недоступна
(таблица~\ref{tab:rentbuy}). Малой команде вместо целого узла B300 подходят
серверы на RTX PRO 6000. Чисто серверный режим на таких серверах мы считаем для закрытого контура, где вся работа идёт на
открытой модели (абзац о закрытом контуре ниже). Команде из пяти человек нужна одна такая карта при лёгкой нагрузке и три при тяжёлой, и свой сервер обходится ей в \$147--437 на человека в месяц против \$3\,373 на человека при лёгкой нагрузке с одним узлом B300.
Команде из двадцати человек нужна одна карта при лёгкой нагрузке и десять при тяжёлой, и свой
сервер обходится ей в \$39--367 на человека в месяц против \$845 на человека при лёгкой нагрузке с одним узлом B300.

\begin{figure}[t]
\centering
\begin{tikzpicture}
\begin{axis}[name=p0,width=0.32\textwidth,height=50mm,xmode=log,ymode=log,ytick={20,50,100,200,500},yticklabels={20,50,100,200,500},xmin=4,xmax=6500,xtick={5,50,500,5000},xticklabels={5,50,500,5000},xlabel={человек},title={200 вызовов в день},tick label style={font=\footnotesize},label style={font=\footnotesize},title style={font=\footnotesize},legend style={font=\footnotesize,draw=none},legend pos=south west]
\addplot[black,mark=*,mark size=1pt] table[x=team,y=light_own,col sep=comma] {data/team-perimeter.csv};
\addplot[black,dashed,thick] table[x=team,y=light_hybrid,col sep=comma] {data/team-perimeter.csv};
\legend{свой кластер,гибрид}
\end{axis}
\begin{axis}[at={(p0.east)},anchor=west,xshift=3mm,name=p1,width=0.32\textwidth,height=50mm,xmode=log,ymode=log,ytick={20,50,100,200,500},yticklabels={20,50,100,200,500},xmin=4,xmax=6500,xtick={5,50,500,5000},xticklabels={5,50,500,5000},xlabel={человек},title={800 вызовов в день},tick label style={font=\footnotesize},label style={font=\footnotesize},title style={font=\footnotesize},legend style={font=\footnotesize,draw=none},legend pos=south west]
\addplot[black,mark=*,mark size=1pt] table[x=team,y=medium_own,col sep=comma] {data/team-perimeter.csv};
\addplot[black,dashed,thick] table[x=team,y=medium_hybrid,col sep=comma] {data/team-perimeter.csv};
\end{axis}
\begin{axis}[at={(p1.east)},anchor=west,xshift=3mm,name=p2,width=0.32\textwidth,height=50mm,xmode=log,ymode=log,ytick={20,50,100,200,500},yticklabels={20,50,100,200,500},xmin=4,xmax=6500,xtick={5,50,500,5000},xticklabels={5,50,500,5000},xlabel={человек},title={2500 вызовов в день},tick label style={font=\footnotesize},label style={font=\footnotesize},title style={font=\footnotesize},legend style={font=\footnotesize,draw=none},legend pos=south west]
\addplot[black,mark=*,mark size=1pt] table[x=team,y=heavy_own,col sep=comma] {data/team-perimeter.csv};
\addplot[black,dashed,thick] table[x=team,y=heavy_hybrid,col sep=comma] {data/team-perimeter.csv};
\end{axis}
\end{tikzpicture}
\caption{Доллары на человека в месяц в закрытом контуре для всех размеров команды. M5 Max 64~ГБ уже есть; оба маршрута выбирают сервер и хранилище при одинаковых сроках хранения и бюджете задержки, без ограничения команд для RTX PRO 6000. Доля открытой модели 100\%, хранение по политике протокола.}
\label{fig:teams}
\end{figure}

\textbf{Аренда или покупка сервера.} Своё железо нужно гибриду, когда команда держит собственный prefill-сервер,
например ради контроля над данными. Небольшие команды обычно работают удалённо, и сервер им удобнее держать в
колокации. Таблица~\ref{tab:rentbuy} сравнивает почасовую аренду RTX PRO 6000 с покупкой такого же
сервера в колокацию при доле 50\%. При рабочем дне в 10 часов команде из 5--10 человек аренда дешевле покупки при любой из рассмотренных остаточных
стоимостей: даже при 80\% покупка дороже на 7\%. На 50 человек аренда (\$33) почти равна покупке только при
остаточной стоимости 80\% (\$35); при 50 и 20\% покупка дороже в 1.5 и 1.9 раза. Если та же работа распределена на
круглые сутки, например между часовыми поясами, хватает меньшего числа серверов, и покупка выигрывает при любом
размере команды.

\begin{table}[t]
\centering
\caption{Серверная часть гибрида на человека в месяц при доле 50\%, тяжёлый профиль. Аренда RTX PRO 6000 по
\$2.20 в час; покупка сервера с одной RTX PRO 6000 в колокацию, амортизация за четыре года с остаточной стоимостью
20, 50 или 80\%. В строках с 24 ч та же дневная нагрузка, распределённая на сутки в рабочие дни. В скобках число серверов.
Хранение копий сессий по политике протокола входит в обе колонки.}
\label{tab:rentbuy}
\footnotesize
\begin{tabular}{lrrrrr}
\toprule
Команда, нагрузка & Аренда & Покупка, 20\% & Покупка, 50\% & Покупка, 80\% \\
\midrule
5 человек, 10 ч в день & \$98 (1) & \$194 (1) & \$149 (1) & \$104 (1) \\
5 человек, 24 ч & \$231 (1) & \$198 (1) & \$153 (1) & \$108 (1) \\
10 человек, 10 ч в день & \$51 (1) & \$100 (1) & \$77 (1) & \$55 (1) \\
10 человек, 24 ч & \$120 (1) & \$104 (1) & \$81 (1) & \$59 (1) \\
50 человек, 10 ч в день & \$33 (3) & \$62 (3) & \$48 (3) & \$35 (3) \\
50 человек, 24 ч & \$53 (2) & \$47 (2) & \$38 (2) & \$29 (2) \\
\bottomrule
\end{tabular}
\end{table}

\textbf{Закрытый контур.}

\textbf{Архивный расчёт закрытого контура с Q4.} Часть организаций не может отправлять данные за свой периметр: ни облачным API, ни
стороннему поставщику prefill. Для них доля 100\% становится обычным случаем, и сравнивать приходится три варианта
внутри периметра: одни клиентские машины, свой серверный кластер и гибрид со своими prefill-узлами. Вызовы, которые
клиент не успевает сгенерировать, гибрид отправляет на генерацию своему же серверу. Таблица~\ref{tab:perimeter}
показывает варианты для трёх профилей нагрузки и клиента M5 Max 64~ГБ, который у сотрудника уже есть.

Одни клиентские машины почти ничего не стоят, но пользователю с 800 вызовами в день
не отвечают на $800-780.7478=19.2522$ вызова, то есть $19.2522/800=2.41\%$;
обслужено 97.59\%, округлённо 98\%, как в таблице. Клиентский prefill занимает в среднем 11.4 секунды.
Гибрид отвечает на все вызовы. При средней нагрузке расчёт даёт 2.1 секунды среднего ожидания prefill
без очереди. Для средней нагрузки гибрид в 2.3--2.5 раза дешевле своего серверного кластера на 50, 500 и 5000 человек,
при осторожной политике хранения в 2.0--2.2 раза. Для лёгкой нагрузки на 50 и 5000 человек гибрид дешевле кластера на
27--39\%, при осторожной политике на 17--30\%. Серверы RTX PRO 6000 и B300 допустимы без ограничения размера команды,
и движок выбирает самый дешёвый из допустимых вариантов. Обе схемы имеют одинаковые сроки хранения и число сохраняемых сессий,
одинаковые ограничения среднего и p90 времени prefill, загрузку не выше 70\%.
Обе схемы сохраняют четыре сессии на человека полный месяц; бюджет среднего prefill равен 19.9 секунды, p90 54.3 секунды.
Ограничение задержки не отсекло ни одного из 2106 вариантов: сравнение определяется стоимостью при равном хранении.
Пул сохранённых сессий оплачен отдельно; занятость GPU включает активную генерацию.
На 500 человек при лёгкой нагрузке гибрид дешевле своего кластера: свой кластер стоит \$32 на человека в месяц против \$20 у гибрида.
При осторожной политике на 500 человек гибрид стоит \$25 против \$35 у своего кластера.
Экономия на лёгкой нагрузке растёт с размером команды. Целочисленные серверы учитываются в цене.
При тяжёлой нагрузке большая часть вызовов генерируется на сервере, и гибрид дешевле кластера на 25--26\%,
при осторожной политике хранения на 24--26\%. Команде из пяти человек свой сервер обходится
от \$147 на человека в месяц при лёгкой нагрузке.

\begin{table}[t]
\centering
\caption{Закрытый контур: всё внутри периметра, доля открытой модели 100\%, клиент M5 Max 64~ГБ уже есть.
Стоимость на человека в месяц; одни клиенты показаны долей обслуженных вызовов за восьмичасовой день.
Одна-две RTX PRO 6000 или узлы B300 допускаются без ограничения размера команды;
из допустимых вариантов взят самый дешёвый. Обе схемы имеют одинаковые ограничения среднего и p90 времени prefill,
сроки хранения и число сохраняемых сессий; рассматриваются HBM, RAM и NVMe.
Остаточная стоимость 20\%, место в стойке уже есть, загрузка не выше 70\%, очереди исключены.}
\label{tab:perimeter}
\footnotesize
\begin{tabular}{rrrrrr}
\toprule
Вызовов в день & Команда & Свой кластер & Гибрид & Гибрид, осторожно & Одни клиенты: обслужено \\
\midrule
200 & 50 & \$34 & \$24 & \$34 & 100\% \\
200 & 500 & \$32 & \$20 & \$25 & 100\% \\
200 & 5000 & \$32 & \$20 & \$25 & 100\% \\
800 & 50 & \$120 & \$53 & \$63 & 98\% \\
800 & 500 & \$116 & \$47 & \$57 & 98\% \\
800 & 5000 & \$114 & \$46 & \$57 & 98\% \\
2\,500 & 50 & \$359 & \$265 & \$272 & 31\% \\
2\,500 & 500 & \$349 & \$260 & \$270 & 31\% \\
2\,500 & 5000 & \$347 & \$259 & \$270 & 31\% \\
\bottomrule
\end{tabular}
\end{table}

\textbf{Цена безубыточности prefill.} На измеренных клиентах при доле 20\% гибрид блокируют накладные расходы даже при бесплатном prefill.
Таблица~\ref{tab:breakeven} раскладывает нулевую цену вычисления при неизменной маршрутизации.
Отсутствие порога обозначено как no break-even; стоимость передачи, хранения, обслуживания, миграции и электричества уже превышает API.
У M5 Max при доле 50\% остаётся порог \$0.226 за миллион токенов, ниже всех измеренных серверных затрат.
Против медианного API при доле 20\% порог равен \$5.0--5.4 за миллион токенов для уже имеющегося клиента при хранении по политике протокола.
Серверная стоимость миллиона токенов prefill по измеренным кривым: llama.cpp PRO 6000 \$0.62--0.63, vLLM PRO 6000 \$0.26, vLLM B300 \$0.31, vLLM B300 с кэшем FP8 \$0.23.
Цены vLLM относятся только к серверу (server-only); мост не квалифицирован по времени полного вызова.
Измеренные кривые взвешены новыми токенами и ответами клиента, с экстраполяцией за измеренную глубину и тарифом при загрузке 70\%.
У оценочной Ultra 256~ГБ при доле 20\% порог против самого дешёвого API равен \$0.116; вывод о блокирующих накладных расходах на эту строку не распространяется.

\begin{table}[t]
\centering
\caption{Месячные доллары при бесплатном prefill. ‡ M3: промежуточные многопоточные опоры PRI~4, повтор при PRI~31 ожидается. M5: приоритет не записан; запускающий процесс показал PRI~31 при последующей проверке. Клиент уже есть, хранение по политике протокола. Столбцы: M3 Max и M5 Max 64~ГБ при долях 20 и 50\%. Вызовы сверх ёмкости оплачены у того же самого дешёвого API. Нет порога означает no break-even.}
\label{tab:breakeven}
\scriptsize
\begin{tabular}{lrrrr}
\toprule
Компонент & M3, 20\% & M5, 20\% & M3, 50\% & M5, 50\% \\
\midrule
Передача & \$4.41 & \$3.23 & \$5.59 & \$6.52 \\
Хранение & \$4.56 & \$4.51 & \$4.56 & \$4.51 \\
Обслуживание & \$5.00 & \$5.00 & \$5.00 & \$5.00 \\
Миграция & \$1.62 & \$1.06 & \$2.04 & \$2.10 \\
Электричество клиента & \$2.52 & \$2.52 & \$2.52 & \$2.52 \\
Вызовы сверх ёмкости, API & \$0.00 & \$0.12 & \$17.35 & \$7.04 \\
Всего при бесплатном prefill & \$18.11 & \$16.44 & \$37.07 & \$27.70 \\
Самый дешёвый API & \$14.10 & \$14.10 & \$35.24 & \$35.24 \\
Порог, долларов за млн токенов & нет порога & нет порога & нет порога & \$0.226 \\
\bottomrule
\end{tabular}
\end{table}

\textbf{Как это будет меняться.} Облачный пользователь получает прогресс моделей в виде снижения цены: по оценке
Epoch AI стоимость фиксированного уровня способностей с 2023 года падает примерно в 13 раз в год
\cite{epoch2026price}. Пользователь гибрида получает его в виде новых способностей на том же железе: модель класса 27B,
которая помещается на тот же Мак, за полгода выросла по индексу Artificial Analysis с 34.6 до 52.0 \cite{aaindex}
(рисунок~\ref{fig:progress}). Цены API снижаются, и уже сегодня самый дешёвый поставщик обходится дешевле гибрида на всех
рассмотренных клиентах. Открытые модели умнеют, доля открытой модели растёт, и с ней растёт экономия против медианного
поставщика. Контекст становится глубже, и налог на поддержание кэша растёт вместе с глубиной, а работа гибридного
сервера растёт медленнее: он заново читает только новые токены и ответы клиента, хотя цена каждого нового токена тоже
растёт с глубиной. Поэтому главный рычаг гибрида находится на глубоком контексте
(раздел~\ref{sec:million}).

\begin{figure}[t]
\centering
\begin{tikzpicture}
\begin{axis}[
  name=price, width=0.47\textwidth, height=48mm, ymode=log, xmin=0, xmax=8.5, ymin=0.1, ymax=1.2,
  xtick={0,2,4,6,8}, xticklabels={фев,апр,июн,авг,окт}, ytick={0.1,0.2,0.5,1}, yticklabels={0.1,0.2,0.5,1},
  ylabel={цена, доля от февраля}, title={API: цена фиксированного уровня},
  tick label style={font=\footnotesize}, label style={font=\footnotesize}, title style={font=\footnotesize}]
\addplot[black,dashed,thick,domain=0:8,samples=40] {13^(-x/12)};
\end{axis}
\begin{axis}[
  at={(price.east)}, anchor=west, xshift=14mm, width=0.47\textwidth, height=48mm, xmin=0, xmax=8.5, ymin=30, ymax=56,
  xtick={0,2,4,6,8}, xticklabels={фев,апр,июн,авг,окт}, ylabel={индекс AA v4.1.1},
  title={Гибрид: модель 27B на том же Маке},
  tick label style={font=\footnotesize}, label style={font=\footnotesize}, title style={font=\footnotesize}]
\addplot[black,thick,mark=*,mark size=1.6pt] coordinates {(0.75,34.6) (2.7,37.7) (6.45,52.0)};
\node[font=\scriptsize,anchor=north west] at (axis cs:0.75,34.6) {Qwen3.5};
\node[font=\scriptsize,anchor=south east] at (axis cs:2.7,37.7) {Qwen3.6};
\node[font=\scriptsize,anchor=south east] at (axis cs:6.45,52.0) {Qwen3.8};
\end{axis}
\end{tikzpicture}
\caption{Как пользователь получает прогресс моделей в 2026 году. Слева: цена фиксированного уровня способностей через
API при снижении в 13 раз в год \cite{epoch2026price}, нормирована на февраль. Справа: индекс Artificial Analysis v4.1.1
моделей класса 27B, которые помещаются в одну и ту же машину с 32~ГБ памяти \cite{aaindex}; значение для Qwen3.5-27B
оценочное.}
\label{fig:progress}
\end{figure}

\textbf{Чувствительность к формату передачи.} Таблица~\ref{tab:bf16-economics} меняет только рекуррентные тензоры
на проводе и пересчитывает политику выбора места чтения, догоняющие вызовы, ёмкость, трафик и цены.
Резидентное состояние, генерация и закон вставки сохраняются. Базовые расчёты выше используют округлённые 157~МБ
состояния; парное сравнение ниже использует точные тензоры на проводе при той же оценке резидентной памяти.
Базовая серверная экономика остаётся на измеренных кривых B300 и PRO 6000;
пороги таблицы~\ref{tab:model-link} отдельно используют одну сборку PRO 6000.
Эти расчёты не измеряют качество и скорость передачи bf16 на GPU или Mac.

\begin{table}[t]
\centering
\caption{Чувствительность экономики к f32/bf16 на проводе. M3 Max, четыре потока, 8 часов, 1~Гбит/с.
Все значения промежуточные: многопоточные опоры M3 сняты при PRI~4; повтор при PRI~31 ожидается.
Трафик относится к тяжёлому профилю; цены для команды из пяти человек, доли открытой модели 20\% и уже купленного Mac.
Базовый тариф трафика \$5/ТБ; строка \$50/ТБ показана отдельно. Модель сохраняет скорость вставки и генерации;
качество и скорость bf16 на GPU/Mac ещё не измерены.}
\label{tab:bf16-economics}
\footnotesize
\begin{tabular}{lrr}
\toprule
Величина & f32 & bf16 \\
\midrule
Локальных сессий в день & 3.972504 & 3.972504 \\
Гибридных сессий в день & 5.469343 & 5.475793 \\
Отношение ёмкости дня & 1.376800 & 1.378424 \\
Ёмкость вызовов в день / обслужено в тяжёлом профиле & 634.443780 & 635.191961 \\
Трафик, ТБ в месяц & 1.118012 & 0.808498 \\
Трафик по тарифу \$50/ТБ, долларов & 55.900580 & 40.424900 \\
Доля оффлоадов & 0.314480 & 0.328463 \\
Трафик по базовому тарифу, долларов & 4.405479 & 3.182101 \\
Хранение и миграция, долларов & 6.182218 & 6.232401 \\
Гибрид всего, долларов & 73.409243 & 72.236048 \\
\bottomrule
\end{tabular}
\end{table}

При bf16 оффлоад выбирается чаще, поэтому стоимость хранения и миграции растёт;
месячный объём передачи уменьшается менее чем вдвое. Экономия байтов на одном оффлоаде
не задаёт отношение месячного трафика при изменившемся числе удалённых вызовов.

% checker-site: v5a-R16-conclusion
При доле открытой модели 20\% гибрид на машине, которая уже есть, в 1.2--1.5 раза дешевле медианного поставщика
API той же модели и в 4.8--6.0 раза дороже самого дешёвого, в зависимости от класса клиента и политики хранения. Эти соотношения задают границы экономической применимости. На глубоком контексте: сессия до миллиона
токенов в команде от 500 человек со своими узлами prefill при полной загрузке Studio обходится в 2.6--2.7 раза дешевле API и
в 1.3 раза дешевле своего кластера. Клиенту, который потерял и память, и диск, сервер возвращает сессию в
миллион токенов примерно за пять минут и дольше, а M3 Max читал бы такую сессию Qwen3.8-27B заново около 20 часов. В закрытом контуре гибрид при средней нагрузке в 2.0--2.5 раза дешевле своего серверного кластера. При лёгкой нагрузке 200 вызовов в день на человека
гибрид на уже купленном M5 Max с 64~ГБ дороже своего кластера на 5.1\% для 5 человек, на 10.1\% для 10 и на 19.5\% для 20
при обеих политиках хранения; для 50 человек гибрид дешевле на 27.3\% по политике протокола и на 17.0\% при осторожной,
для 500 дешевле на 38.2\% и 28.9\%, для 5000 дешевле на 38.5\% и 29.8\% соответственно.
Серверы на RTX PRO 6000 допустимы без ограничения размера команды; обе схемы имеют одинаковые требования хранения и задержки. Условия применимости включают уже купленную машину класса Max или Ultra,
крупную модель, глубокий контекст, большие порции от инструментов и канал, на котором передача окупается по времени. Покупать дорогой клиент только ради экономии против самого
дешёвого API невыгодно; доплату к рабочей машине могут оправдать память, полоса и работа без поставщика.

% checker-site: v5a-old-economics-summary
\textbf{Вывод.} Против медианного поставщика API гибрид на клиентах класса Max,
которые уже есть, при доле 20\% в 1.2--1.4 раза дешевле, при доле 50\% в 1.6--2.4 раза.
В последнем случае клиент обслуживает от 51 до 100\% вызовов, остальные уходят к самому дешёвому API.
Для Ultra с хранением по обеим политикам выигрыш против медианного поставщика в 1.4--1.7 раза.
Самый дешёвый поставщик API дешевле гибрида на всех рассмотренных клиентах и при обеих долях.
В закрытом контуре при средней нагрузке гибрид в 2.0--2.5 раза дешевле своего кластера, при тяжёлой на 25--26\%,
при осторожной политике на 24--26\%, а при лёгкой на 50 и 5000 человек гибрид дешевле своего кластера; на 500 человек гибрид тоже дешевле:
серверы на RTX PRO 6000 допустимы без ограничения размера команды, обе схемы имеют одинаковые требования хранения и задержки.

% checker-site: v5a-old-abstract-economics
Экономика на профиле из 1.5 миллиона вызовов задаёт условия применимости: при доле открытой модели 20\%
гибрид на уже купленной машине Max или Ultra в 1.2--1.5 раза дешевле медианного поставщика API той же модели
и в 4.8--6.0 раза дороже самого дешёвого, в зависимости от клиента и политики хранения.

% checker-site: v5a-old-introduction-economics
При доле открытой модели 20\% гибрид на уже купленной машине Max или Ultra в 1.2--1.5 раза дешевле медианного поставщика API той же модели
и в 4.8--6.0 раза дороже самого дешёвого; диапазоны учитывают клиента и политику хранения.
Эти соотношения задают границы экономической применимости.

\textbf{Уточнение времени B300 в введении.} Интерполяция времени на токен по опорам FACTS 1b
1284 токена/с при 100 тысячах, 1131 при 150 тысячах и 1010 при 200 тысячах использует
$t(d)=10000[(1-w)/r_a+w/r_b]$, $w=(d-a)/(b-a)$.
При $d=131000$, $a=100000$, $b=150000$ получаем $w=0.62$ и
$t=10000(0.38/1284+0.62/1131)=8.44$ секунды.
При $d=181000$, $a=150000$, $b=200000$ получаем $w=0.62$ и
$t=10000(0.38/1131+0.62/1010)=9.50$ секунды.
Оба времени рассчитаны по измеренным опорам llama.cpp, отдельными замерами не являются.

\textbf{Архивный M5 Max, комплект 1.4.1.} M5 Max с 64~ГБ и 40-ядерным GPU измерен на macOS 27.0.1 комплектом 1.4.1.
На тех же глубинах сегментный prefill дал 586, 269, 151, 95, 70 и 55 токенов/с,
генерация 23.1, 19.8, 14.6, 12.2, 10.6 и 9.25 токена/с. Добавочный prefill
порции в 2000 токенов снизился с 394 до 51 токена/с между концами этого диапазона.
Запуск со swap объёмом 5.1~ГБ исключён из оценки скорости машины. Приоритет исторического запуска M5 Max не записан.
Логи не позволяют установить, использовались ли тензорные блоки M5.

\textbf{Архивный параллельный M5, комплект 1.6.6.} Приоритет этих опор не записан;
запускающий процесс показал PRI~31 при последующей проверке. Текущие измерения 1.6.7 заменяют эти опоры;
производные расчёты run 16 сохраняют исходные данные до пересчёта.
\begin{center}
\begin{tabular}{llrrrrrrrrr}
\toprule
Клиент, запуск & Глубина & 1 & 2 & 4 & 6 & 8 & 12 & 16 & 24 & 32 \\
\midrule
M5 Max, 1.6.6 & 32 & 22.62 & 31.89 & 47.78 &  & 62.32 & 62.78 & 65.54 &  &  \\
M5 Max, 1.6.6 & 64 & 19.88 & 26.00 & 32.55 &  & 40.23 &  &  &  &  \\
M5 Max, 1.6.6 & 128 & 15.73 & 18.35 & 20.69 &  &  &  &  &  &  \\
\bottomrule
\end{tabular}
\label{tab:m5-sweep-history}
\end{center}

% checker-site: v5a3-old-capacity-conclusion
Мы исследовали устройство и границы применимости prefill как услуги для автономного локального агента.
В рассматриваемой схеме сервер читает новые порции контекста, а машина
пользователя генерирует ответы и хранит память модели. В модельных расчётах сервер присылает клиенту дельту памяти, и большая часть ожидания
prefill уходит из агентского цикла: порцию, которую клиент читал бы десятки секунд, сервер читает за секунды. Клиенту остаётся
генерация, которую ограничивает полоса памяти, и параллельные агенты используют её лучше. Машины с
унифицированной памятью класса Max и Ultra работают в обоих режимах: без сервера автономно, с сервером с меньшим
ожиданием чтения контекста. Схема даёт три разных
результата. Ожидание prefill на M3 Max сокращается на 10--15 секунд на вызов, число вызовов, которые клиент успевает
за день, растёт умеренно, на M3 Max расчётно в 1.38 раза при канале 125~МБ/с (1~Гбит/с)\textsuperscript{‡} (промежуточные многопоточные опоры PRI~4; повтор при PRI~31 ожидается), на M5 Max с 64~ГБ на 29\%, на M5 Max с 128~ГБ по переносу измерений на 28\%; приоритет M5 не записан, запускающий процесс показал PRI~31 при последующей проверке; у Ultra прибавки нет, а экономия денег зависит от того, какую долю работы
можно отдать открытой модели.

\subsection{Архивный сценарий миллиона токенов, run 14}
\label{app:million-history}

\textbf{База расчётов раздела.} Экономика GLM, цены сессий, размеры команд, загрузка Studio,
отношения цен и рисунок~\ref{fig:million} сохраняют прежнюю базу llama.cpp Q4, тарифы API и снимок run 14.
Расчёты описывают архивный сценарий; текущие границы применимости после перехода на vLLM ожидают движка.
Новые измерения восстановления Qwen приведены отдельно от этой экономики.

Чем глубже контекст, тем больше агент держит в голове: весь репозиторий, длинную историю работы, большой набор
документов. Налог на поддержание кэша растёт вместе с глубиной, потому что каждый вызов перечитывает весь контекст. У
поставщика GLM-5.3-Flash сессия, которая дорастает до 200 тысяч токенов, стоит \$0.44, а сессия до миллиона \$9.15:
контекст глубже в пять раз, сессия дороже в 21 раз.

\textbf{Клиент для миллиона токенов.} GLM-5.3-Flash устроена как MoE: 328 миллиардов параметров, из них 17.5 активных
\cite{glm53flash}. Её кэш внимания занимает около 6~КБ на токен, в шесть раз меньше, чем у Qwen3.8-27B, а рекуррентное состояние около 153~МБ в памяти и 76~МБ при передаче в 16 битах. Веса в 4 битах занимают 181~ГБ и помещаются в Mac Studio с M5 Ultra и 256~ГБ памяти. Рядом с весами
помещается 10 сессий по миллиону токенов или 46 сессий по 200 тысяч. Генерация, по нашей оценке, идёт со скоростью
около 19 токенов в секунду на поток при четырёх потоках. Узкое место то же, что и у Qwen3.8-27B: прочитать миллион токенов сама Studio может по нашей оценке примерно за 42 минуты, а серверная часть на B300 вместе с передачей по гигабитному каналу занимает по оценке 454 секунды.
Эта оценка холодной загрузки исключает вставку состояния; перенос скорости B300 на GLM оценочный. Qwen3.8-27B
на клиенте с 128~ГБ по объёму памяти держит две сессии по миллиону токенов, но это за пределами её штатного
контекста (таблица~\ref{tab:fit}).

\textbf{Потеря кэша.} Клиент, который потерял и память, и сохранённое на диске состояние сессии в миллион токенов,
почти теряет саму сессию. M3 Max читал бы такую сессию Qwen3.8-27B заново 20 часов: время на токен по замеренным
отрезкам prefill (таблица~\ref{tab:m3max}) мы продолжаем линейно за 230.9 тысячи токенов с наклоном отрезка от 181.4 до 230.9 тысячи (19.9 часа) или от 131.4 до 230.9 тысячи (19.9 часа). Это оценка за пределами замеров и штатного контекста
модели. В гибриде сервер хранит копию сессии и присылает клиенту её память целиком: на миллионе токенов это около 35~ГБ, по гигабитному каналу около пяти минут, трафик по \$5 за терабайт стоит около 18 центов. Хранение копии на NVMe по \$0.10 за гигабайт стоит около \$3.5 в месяц. По политике протокола копия в рабочие часы лежит в оперативной памяти по \$0.50 за гигабайт, и хранение стоит около \$6.7. Пять минут здесь нижняя граница: по политике протокола
сервер дописывает в копию токены вызовов, прочитанных клиентом, только при следующем оффлоаде, и перед
восстановлением ему может понадобиться prefill этого хвоста. Если диск клиента уцелел, сессию быстрее поднять из
локального снимка на его SSD.

\textbf{Стоимость сессии.} Следующий расчёт сохраняет базу llama.cpp Q4 и прежние тарифы API; текущие цены vLLM ожидают пересчёта. Рисунок~\ref{fig:million} сравнивает стоимость сессии до миллиона токенов у API, у своего
серверного кластера и у гибрида с Studio. Гибрид включает покупку Studio целиком: \$10\,799 с амортизацией за четыре
года при остаточной стоимости 35\% и электричеством дают \$149 в месяц. При двух таких сессиях в день на человека это \$3.54 на сессию только за Studio. Для команд от
500 человек гибрид обходится в \$5.7--6.1 за сессию против \$9.15 у API.
На 500 человек свой кластер стоит \$4.81 за сессию, на 5000 человек \$4.49 и дешевле гибрида.
Команде из 10 человек и меньше выделенный сервер не окупается, и дешевле всего остаётся API. На 50 человек свой кластер
(\$8.02) немного дешевле API (\$9.15), а гибрид с отдельной Studio на каждого дороже обоих (\$11.74). Малым командам нужен общий
поставщик prefill.

Studio при двух сессиях в день загружена на 38\%: в гибридном режиме она успевает провести почти шесть сессий по миллиону токенов за рабочий день. При полной загрузке доля Studio падает до \$1.33 на сессию, и сессия
гибрида для команды от 500 человек стоит \$3.4--3.5, в 2.6--2.7 раза дешевле API и в 1.3 раза дешевле своего кластера. Четыре сессии в день на человека дают \$3.9, в 1.2 раза дешевле своего кластера и в 2.4 раза дешевле API; если Studio делят двое и каждый проводит по две сессии в день, \$4.0--4.4. При осторожной политике хранения сессии дороже не больше чем на 20\%. Загрузка клиента и есть главный
рычаг гибрида на глубоком контексте.

\begin{figure}[t]
\centering
\begin{tikzpicture}
\begin{axis}[
  width=0.62\textwidth, height=55mm, xmode=log, ymode=log,
  xmin=4, xmax=6500, ymin=2, ymax=120, xtick={5,10,20,50,500,5000}, xticklabels={5,10,20,50,500,5000},
  ytick={2,5,10,20,50,100}, yticklabels={2,5,10,20,50,100},
  xlabel={человек в команде}, ylabel={\$ за сессию},
  tick label style={font=\footnotesize}, label style={font=\footnotesize},
  legend pos=north east, legend style={font=\footnotesize,draw=none,fill=none}]
\addplot[black,dashed] table[x=team,y=api,col sep=comma] {data/glm-1m-cost.csv};
\addlegendentry{API}
\addplot[black,mark=*,mark size=1.4pt] table[x=team,y=own,col sep=comma] {data/glm-1m-cost.csv};
\addlegendentry{свой кластер B300}
\addplot[black,very thick] table[x=team,y=hybrid,col sep=comma] {data/glm-1m-cost.csv};
\addlegendentry{гибрид, Studio куплена}
\end{axis}
\end{tikzpicture}
\caption{Архивный сценарий run 14, llama.cpp Q4 и прежние тарифы API; текущая экономика ожидает пересчёта. Стоимость сессии GLM-5.3-Flash, которая дорастает до миллиона токенов, в зависимости от размера команды. Две
такие сессии в день на человека; гибрид включает покупку Mac Studio M5 Ultra 256~ГБ и серверный prefill на своих узлах
B300.}
\label{fig:million}
\end{figure}

\textbf{Вывод архивного сценария.} Следующие отношения относятся к run 14, llama.cpp Q4 и прежним тарифам API; текущую экономику vLLM эти отношения не задают. Контекст до миллиона токенов в этом расчёте выгоден гибриду, когда клиент загружен: при четырёх сессиях в день на
Studio гибрид в 1.2 раза дешевле своего кластера и в 2.4 раза дешевле API, при полной загрузке в 1.3 и
в 2.6--2.7 раза.\footnote{Снимок run 14: \texttt{glm\_utilization.csv}, контекст 1000000, команды 500 и 5000, одна Studio на пользователя, четыре сессии в день; отношение полей \texttt{api\_session/hybrid\_session}.} При двух сессиях в день гибрид дороже своего кластера на 27\%, при одной в 2.1--2.2 раза;
малым командам без общего поставщика prefill выгоднее API.

\label{v5a4-m4-row}
\textbf{Исключённая оценка M4 Max из run 16.} Строка сохраняется только для истории расчёта;
8.38 сессии без сервера у оценённого M4 против 6.73 у M5 на измеренных опорах требует пересмотра оценки prefill движком.
\begin{center}\begin{tabular}{lrrrr}
M4 Max 128 ГБ & 19.5 & 8.38 & 11.01 & 1.31 \\
\end{tabular}\end{center}

\subsection{Архивная синтетическая медиана тарифов}
% checker-site: v5b-old-synthetic-median
Следующий абзац сохраняет прежнюю синтетическую медиану отдельных полей тарифа и прежний профиль.
Текущая экономика выбирает реальных допущенных поставщиков отдельно в каждой точке кэша.

Облачный API снимает ожидание prefill. На B300 измеренная кривая llama.cpp даёт для порции
10 тысяч токенов расчётные 8.44 секунды на глубине 131 тысяча и 9.50 секунды на 181 тысяче; vLLM FP8 показывает другие
серверные скорости: на глубинах профиля отношение vLLM FP8 к llama.cpp Q4 составляет расчётные 10.1--10.8 раза по измеренным опорам (раздел~\ref{sec:server}). Поставщик
хранит кэш контекста между вызовами, но берёт плату за его чтение на каждом вызове. В архивной базе llama.cpp Q4 плата за чтение кэша составляет большую часть расходов агентской работы на глубоком контексте у медианного поставщика. На нашем профиле средняя глубина контекста 164 тысячи токенов, а новых токенов в вызове около тысячи. Даже со
скидкой на чтение кэша вызов перечитывает примерно в 150 раз больше токенов, чем добавляет: у синтетического медианного поставщика Qwen3.8-27B, составленного из медианы каждого ценового поля снимка от 29 октября 2026 года (\$0.05 за миллион токенов чтения кэша,
\$0.18 за миллион новых и \$2.30 за миллион выходных) чтение кэша даёт около 85\% цены вызова, у самого дешёвого около 48\%. Обычный API не раскрывает промежуточное состояние, которое сервер использовал для генерации.
Клиент с открытыми весами может повторить вычисление; получение состояния даёт объект для отдельного аудита.

\subsection{Архивная ёмкость клиентского дня}
\label{app:client-capacity-history}

\textbf{Скорость генерации и её запас.} На каждом шаге генерация читает веса, которые участвуют в шаге, и весь кэш
контекста, поэтому её предел задаёт полоса памяти. У Qwen3.8-27B это все веса, около 19~ГБ. На M3 Max с полосой 400~ГБ/с средняя генерация при PRI~31 дала 14.99 токена/с на глубине 16 тысяч. Несколько потоков читают веса один раз на всех, поэтому суммарный
предел растёт с числом потоков: для четырёх потоков на глубине 32 тысячи он около 64 токенов в секунду. Охлаждённый замер при PRI~4 дал 24.03 токена в секунду на четырёх потоках против 8.77 у одного,
около 37\% этого предела. Насколько движки сократят разрыв между замером и пределом, покажут следующие замеры;
расчёты используют измеренную скорость генерации. Когда чтение порций уходит на сервер, работа клиента
сводится к генерации, и ускорение генерации движком сокращает время всей схемы.

\textbf{Сколько успевает клиент за день.} Архивная модель run 16 сохраняет прежние опоры M5; новые измерения 1.6.7 ожидают пересчёта движка. Все оценки M3\textsuperscript{‡} в таблицах~\ref{tab:clients} и \ref{tab:agents-day}, включая доли зон и отношения ёмкости, промежуточные: многопоточные опоры сняты при PRI~4; ограниченный повтор при обычном приоритете измерен; обновление архивной модели ожидается. Таблица~\ref{tab:clients} показывает модельную ёмкость
восьмичасового дня с четырьмя агентами, включая вставку состояния. На M3 Max локальная ёмкость
равна 3.97 условной сессии, гибридная 5.47, расчётный выигрыш 1.38 раза при канале 125~МБ/с (1~Гбит/с). Средняя экономия prefill на всём профиле равна 10--15 секундам на вызов в зависимости от канала и формата состояния
(таблица~\ref{tab:offload}). M5 Max с 128~ГБ использует измеренные кривые той же SoC M5 Max с 64~ГБ;
M5 Ultra масштабирует эти кривые по вычислениям и полосе памяти. Это перенос измерений между
конфигурациями, без отдельного замера M5 Max 128~ГБ или Ultra. Отношение сегментного prefill архивного M5 Max комплекта 1.4.1 с незаписанным приоритетом к M3 Max при обычном приоритете
на сопоставимых глубинах лежит между $95.1/52.5=1.81$ и $586.4/191.5=3.06$;
участие тензорных блоков M5 логи не показывают. В архивной модели run 16 на гигабите M5 Max получает 28--29\% к ёмкости дня, а M5 Ultra с учётом вставки не получает прибавки.
NVIDIA DGX Spark на GB10 Grace Blackwell имеет 128~ГБ унифицированной памяти LPDDR5X
и полосу 273~ГБ/с \cite{nvidiadgxspark}; скорости Spark и Strix Halo остаются оценками.
Ноутбуки RTX Spark на N1X продаются с 16 октября 2026 года \cite{nvidiartxsparknews},
до 128~ГБ памяти \cite{nvidiartxspark}. Полоса памяти RTX Spark не опубликована;
расчёт RTX Spark отсутствует.

\begin{table}[t]
\centering
\caption{Условных сессий по 116 вызовов (среднее число вызовов в сессии в журналах) за восьмичасовой день с четырьмя
параллельными агентами, Qwen3.8-27B, канал 1~Гбит/с.
Генерация одного потока на глубине 32 тысячи. M3 использует измеренную зависимость от глубины
и числа потоков, M5 использует измеренный M5 Max с указанным переносом на Ultra.
Вставка состояния включена; коэффициент M3 Max перенесён на остальные клиенты как некалиброванная оценка.
Ёмкость дня и выигрыш сохраняют архивную модель run 16; новые скорости 1.6.7 ожидают пересчёта движка. ‡ M3: промежуточные оценки с многопоточными опорами PRI~4, ограниченный повтор при обычном приоритете измерен; обновление архивной модели ожидается. Архивная модель M5 использует прежние опоры без записи приоритета; новые измерения M5 при PRI~31 приведены отдельно.}
\label{tab:clients}
\footnotesize
\begin{tabular}{lrrrr}
\toprule
Клиент & Генерация, ток/с & Сессий сам & Сессий в гибриде & Выигрыш \\
\midrule
M3 Max 128 ГБ & 14.3 & 3.97 & 5.47 & 1.38 \\
M5 Max 64 ГБ & 22.6 & 6.73 & 8.67 & 1.29 \\
M5 Max 128 ГБ & 22.6 & 6.58 & 8.44 & 1.28 \\
M5 Ultra 96 ГБ & 45.2 & 12.95 & 12.95 & 1.00 \\
M5 Ultra 256 ГБ & 45.2 & 13.64 & 13.64 & 1.00 \\
AMD Strix Halo 128 ГБ & 9.1 & 5.04 & 5.57 & 1.11 \\
NVIDIA DGX Spark & 9.7 & 5.94 & 5.94 & 1.00 \\
\bottomrule
\end{tabular}
\end{table}

\textbf{Выигрыш по числу агентов.} Таблица~\ref{tab:agents-day} использует полную смесь
39\,904 вызовов, по 116 вызовов на условную сессию. Модель учитывает общий ресурс клиента
и время цикла отдельного агента, включая инструменты и вставку состояния.
Строка 16 агентов остаётся модельной: 47.4\% вызовов и 63.9\% времени генерации лежат глубже измеренных опорных точек.
На четырёх агентах эти доли равны 9.3\% и 22.4\%. За последней точкой каждого измеренного числа потоков
модель продолжает скорость по глубинной зависимости; отсутствующее число потоков использует
ближайшее меньшее измеренное. Доли зон пересекаются: при 16 агентах подстановка числа потоков
затрагивает 24.1\% вызовов и 39.7\% времени генерации, ограничение малой глубины первой точкой 3.3\% и 1.2\%.
Измеренная зона покрывает 44.75\% вызовов и 25.73\% времени генерации. Эти доли описывают допущения модели.

\begin{table}[t]
\centering
\caption{M3 Max\textsuperscript{‡}, промежуточные оценки по многопоточным опорам PRI~4; ограниченный повтор при обычном приоритете измерен; обновление архивной модели ожидается. Полная смесь вызовов, восьмичасовой день, канал 1~Гбит/с. Последние два столбца показывают долю зоны моделируемой глубины среди вызовов и среди времени генерации на устройстве. Строка 16 модельная: хотя бы один флаг имеют 55.2\% вызовов и 74.3\% времени генерации. Пересекающиеся флаги строки 16 (вызовы / генерация):
моделируемая глубина 47.4 / 63.9\%, подстановка числа потоков 24.1 / 39.7\%,
нижняя граница 0.0 / 0.0\%, ограничение малой глубины 3.3 / 1.2\%.}
\label{tab:agents-day}
\scriptsize
\begin{tabular}{rrrrrr}
\toprule
Агентов & Сессий сам & В гибриде & Выигрыш & Модель: вызовы & Модель: генерация \\
\midrule
1 & 3.06 & 3.74 & 1.22 & 10.3\% & 18.0\% \\
2 & 3.18 & 4.00 & 1.26 & 9.3\% & 15.2\% \\
4 & 3.97 & 5.47 & 1.38 & 9.3\% & 22.4\% \\
6\textsuperscript{†} & 2.50 & 3.02 & 1.21 & 9.3\% & 12.2\% \\
8 & 4.45 & 6.42 & 1.44 & 9.3\% & 27.5\% \\
12 & 4.73 & 7.01 & 1.48 & 47.4\% & 62.6\% \\
16 & 4.80 & 7.18 & 1.49 & 47.4\% & 63.9\% \\
\bottomrule
\end{tabular}
\par\smallskip
\begin{minipage}{\linewidth}
\noindent\footnotesize † Для шести агентов измеренная зона покрывает 0\% вызовов;
90.67\% вызовов используют удержанную суммарную скорость 6.68 токена/с из точки при 260 тысячах.
Строка сохраняется в таблице и исключена из трендов и отношений роста ёмкости по числу агентов.\normalsize
\end{minipage}
\end{table}

Переход от одного к четырём, восьми и шестнадцати агентам увеличивает расчётную локальную ёмкость
в 1.30, 1.46 и 1.57 раза, гибридную в 1.46, 1.72 и 1.92 раза.
Переход от четырёх к шестнадцати даёт 1.21 и 1.31 раза соответственно. ‡ Эти отношения промежуточные, многопоточные опоры PRI~4; ограниченный повтор при обычном приоритете измерен; обновление архивной модели ожидается.

\subsection{Архивная оценка рабочего дня}
Следующие оценки сохраняют архивную модель run 16 с прежними опорами M5; текущая экономика выделенного сервера дана в разделе~\ref{sec:economics}. Расчёты исходят из того, что оператор участвует в агентском цикле и работает восьмичасовой день. Многие процессы
уже идут без оператора: агенты на автопилоте пишут и проверяют код, собирают и разбирают данные круглые сутки. Такая работа поднимает
загрузку серверов и клиентских машин, и цена на вызов в гибриде от этого падает: за сутки M5 Max обслуживает втрое
больше вызовов, около 3\,020 вместо 1\,010. Автоматическую работу можно
отдавать открытой модели в гибридной схеме, а результат проверять облачной моделью с участием оператора; в наших расчётах этот сценарий не учтён, и в этом отношении они занижают выгоду гибрида. Другие допущения действуют
в обратную сторону. Хранение копий сессий на сервере задано оценкой цены памяти и политикой хранения, которую протокол
допускает, а прототип не реализует; при осторожной политике гибрид дороже на \$10--29 на человека в месяц
(раздел~\ref{sec:economics}).

\subsection{Архивные расчёты B300 для одного хода}
\label{app:channel-model-history}

\textbf{Один ход.} Таблица~\ref{tab:turn} показывает, сколько ждёт агент на одной порции у клиента и через сервер.
Порция в 10 тысяч новых токенов на глубине 200 тысяч занимает у M3 Max 4.9 минуты; через канал в 100~Мбит/с гибрид
отдаёт её за 58 секунд, через гигабит за 22 секунды. Малая порция в 2 тысячи токенов выигрывает меньше, потому что
основную часть дельты составляет рекуррентное состояние, а его размер от порции не зависит.

\begin{table}[t]
\centering
\caption{Время одной порции на M3 Max и через гибрид: подъём копии сессии в память GPU, prefill на одной B300,
передача дельты, вставка и последний токен на клиенте. В скобках ускорение. Qwen3.8-27B Q4, кэш в 8 битах, состояние в 32 битах; время на клиенте по замеренной
кривой prefill.}
\label{tab:turn}
\footnotesize
\begin{tabular}{lrrrrr}
\toprule
Порция, глубина & Клиент & 100 Мбит/с & 300 Мбит/с & 1 Гбит/с & 10 Гбит/с \\
\midrule
2 тыс., 100 тыс. & 30.4 с & 23.7 с (1.3×) & 11.6 с (2.6×) & 7.4 с (4.1×) & 5.7 с (5.3×) \\
10 тыс., 100 тыс. & 162.4 с & 52.8 с (3.1×) & 25.8 с (6.3×) & 16.4 с (9.9×) & 12.7 с (12.8×) \\
2 тыс., 200 тыс. & 54.5 с & 27.7 с (2.0×) & 15.6 с (3.5×) & 11.4 с (4.8×) & 9.7 с (5.6×) \\
10 тыс., 200 тыс. & 295.8 с & 58.5 с (5.1×) & 31.5 с (9.4×) & 22.1 с (13.4×) & 18.5 с (16.0×) \\
\bottomrule
\end{tabular}
\end{table}

\textbf{Разовая загрузка.} Иногда контекст приходит целиком: агент загружает репозиторий или длинный журнал,
запускается новый субагент, сессия начинается заново после сжатия. Тогда клиент читает весь контекст с нуля
(таблица~\ref{tab:oneshot}). 200 тысяч токенов занимают у M3 Max 50 минут, а гибрид через гигабитный канал отдаёт их
за 224 секунды: 157 секунд считает сервер, 10 секунд занимает вставка, остальное приходится на передачу и последний токен. Здесь гибрид выигрывает во времени на любом канале из таблицы; трафика при этом около 7~ГБ на загрузку.

\begin{table}[t]
\centering
\caption{Время до первого токена при загрузке всего контекста сразу, M3 Max и гибрид. Та же модель и те же форматы,
что в таблице~\ref{tab:turn}.}
\label{tab:oneshot}
\footnotesize
\begin{tabular}{lrrrrr}
\toprule
Контекст & Клиент & 100 Мбит/с & 300 Мбит/с & 1 Гбит/с & 10 Гбит/с \\
\midrule
100 тыс. токенов & 14.9 мин & 364 с & 170 с & 102 с & 76 с \\
200 тыс. токенов & 50.4 мин & 737 с & 357 с & 224 с & 173 с \\
262 тыс. токенов & 84.2 мин & 979 с & 484 с & 310 с & 244 с \\
\bottomrule
\end{tabular}
\end{table}

\subsection{Условия сценарных границ и расчёты из введения}
\label{app:v5d6-calculations}
Следующие значения описывают компонент открытой модели при доле 20\%.
В закрытом контуре остальные 80\% обслуживает внутренняя большая модель с одинаковой ценой в обоих вариантах;
расходы этой модели сокращаются при сравнении. Полная доля 100\% и смешанная схема рассчитаны отдельно в подразделе~\ref{sec:closed-contour},
с запретом серверной генерации и явным отказом по ёмкости.
Данные расчётов 22--24 сохранены до проверки парка и покрытия; текущие цены и занятость даны в разделе~\ref{sec:economics}.
\textbf{Численные условия компонента 20\%.} Расчёт относится к команде с уже имеющимися M5 Max 64~ГБ,
собственным выделенным RTX PRO 6000 с vLLM FP8 и предполагаемым каналом 125~МБ/с, RTT 1~мс.
Лёгкая нагрузка составляет 200 вызовов в день, средняя 800; тяжёлая, 2500, соответствует интенсивности владельца.
Лёгкая и средняя нагрузки заданы как сценарии. Доля открытой модели 20\%; сравниваются формы вызовов внутри общей измеренной лестницы,
10\,885 из 39\,904 (27.28\%), с приведением к заданному числу вызовов. Полная исходная смесь этим расчётом не оценена.

Для компонента открытой модели при доле 20\% гибрид дешевле своего кластера на том же сервере
при средней нагрузке с пяти человек, при тяжёлой с двух, в проверенном диапазоне до 500 человек.
При средней нагрузке и пяти людях цена \$130.33 на человека в месяц при рабочей/календарной занятости
6.19/1.42\% против \$245.62 и 32.12/7.39\% у своего кластера. Ответы генерируют Маки, поэтому серверов требуется меньше.
При тяжёлой нагрузке собственный офисный сервер даёт выигрыш против медианного допущенного API
с 45 человек при условной доле кэша 98.1\%; на арендованном сервере с каналом 112.5~МБ/с и RTT 400~мс
такого пересечения до 500 человек нет. Для общего занятого сервера предел распределённых расходов
при доле кэша 98.1\% даёт отношение цены медианного API к гибриду 1.701 при тяжёлой нагрузке
и 0.934 при средней; эти значения используют распределение при занятости 70\%, нижняя оценка при 100\% приведена отдельно.
Сервис общего prefill ведут и авторы статьи; выделенный сервер остаётся основной базой сравнения.

Лёгкая нагрузка оставляет выделенный сервер команды из 50 человек занятым на 3.61\% рабочего времени
и 0.83\% календарного; медианный API при условных 98.1\% кэша примерно в 11 раз дешевле.
Свой кластер на том же сервере также дешевле: обоим маршрутам хватает одинакового парка, гибрид добавляет расходы клиента.
Самый дешёвый допущенный API при условной доле кэша 98.1\% дешевле гибрида примерно в 55 раз при лёгкой нагрузке
и в 4.4 раза при тяжёлой. Без кэша гибрид обходит самый дешёвый допущенный API только при тяжёлой нагрузке
с 74 человек на офисном сервере. При измеренном кэше v1 с долгими паузами гибрид дешевле медианного API при тяжёлой нагрузке с 5 человек;
при лёгкой и средней пересечения до 500 человек нет. В измеренных точках Qwen гибрид не дешевле самого дешёвого API.
Покупка M5 Max добавляет \$69.05 на человека в месяц: пересечение со своим кластером при средней нагрузке
сдвигается с 5 до 34 человек, пересечение с медианным API при условных 98.1\% кэша исчезает.

\textbf{Ожидание и автономность.} Выигрыш M3 Max в статье измерен во времени; строк экономики M3 Max здесь нет.
Порция 10 тысяч токенов на глубине 150 тысяч при обычном приоритете требует около $10000/44.4=225$ секунд.
Передача в bf16 несёт $10000\times34816+78446592=426606592$ байта, около 426.6~МБ:
при 125~МБ/с только передача занимает $426.6/125=3.4$ секунды, при 5~МБ/с $426.6/5=85.3$;
к передаче добавляются серверный prefill и импорт. M5 Max читает быстрее, поэтому на 100 тысячах и порции 2343
с bf16 порог канала равен 34.475~МБ/с против 5.305 у M3 Max при RTT 1~мс.
Измеренный предварительный расчёт на канале 5.19~МБ/с с RTT 125.8~мс оставил все ходы локальными.
На номинальном миллионе токенов расчётный порог M3 Max с bf16 равен 3.351~МБ/с: локальный ход измерен на 1\,040\,000,
а восстановление полного слота с тёплого LaCie служит явно заданной заменой стоимости восстановления, без замера импорта дельты.
Без сервера Мак продолжает работу медленнее; серверная копия сессии позволяет восстановиться после сбоя.

Для владельца ПК с RTX 5090 гибрид мало добавляет при модели 27B и контексте до 262 тысяч:
одна сессия помещается в 32~ГБ, prefill быстрее M3 Max примерно в 23 раза, генерация в 4.3--4.7 раза.
Граница применимости проходит там, где модель и контекст не помещаются в видеопамять: миллион токенов,
несколько сессий одновременно; для крупных MoE это ожидание без замера.
Чужой сервер prefill видит весь контекст пользователя. Данные остаются внутри компании при собственном сервере.


\textbf{Оценка предела генерации на миллионе.} Для M3 Max используем опубликованную полосу
400~ГБ/с и измеренный файл модели 18\,973\,870\,432 байта; это десятичные единицы.
KV q8\_0 при 1\,040\,000 токенах занимает $34816\times1040000=36208640000$ байт,
рекуррентное состояние f32 156\,893\,184 байта. Оценка чтения всех весов и состояния на токен:
$400\cdot10^9/(18973870432+36208640000+156893184)=7.2281$ токена/с.
Измерено 1.6638 токена/с при обычном приоритете: $7.2281/1.6638=4.34$ раза ниже этой оценки.
Ответ на 512 токенов занимает $512/1.6638=307.73$ секунды, около 5.13 минуты.
Это оценка по пиковой полосе без вычислений и накладных расходов, без утверждения об их причинах;
размер файла служит оценкой объёма чтения весов. Полоса указана в спецификации Apple M3 Max \cite{appleM3}.

\textbf{Формула цены узла.} $30000(1-0.2)/(12\times4)+30000\times0.02/12+
0.9\times1.3\times730\times0.5\times0.15=500+50+64.0575=614.0575$ доллара в месяц.
\textbf{Нижняя граница лёгкой нагрузки.} $2.52+5=7.52$ доллара на человека;
$3.604082737\times840/1000=3.03$ доллара у измеренного минимума v1 с одним тёплым наблюдением.

\subsection{Архив неполной RULER до возобновления}
Следующие строки описывают заменённый сохранённый срез до завершения обеих глубин. Текущие результаты 390/390
и смена хоста при завершении 128 тысяч приведены в разделе~\ref{sec:quality}; архив не служит текущим вердиктом.
\textbf{Промежуточные сохранённые строки Qwen на 64 и 128 тысячах.}
Таблица~\ref{tab:ruler-interim-archive} описывает неполные последовательные префиксы задач,
сохранённые до остановки узлов по решению владельца 2026-11-04.
Это промежуточные результаты с числом пар из 390, без зарегистрированного вердикта;
сохранённые строки не являются случайной подвыборкой. В прежнем первичном анализе обе глубины
входили как ожидающие с $p=1$; полный анализ заменяет этот снимок.

\begin{table}[t]
\centering
\caption{Qwen RULER: промежуточные сохранённые пары из 390, отдельно от первичного анализа.
Во всех трёх вариантах оценка и число полностью верных ответов совпали на сохранённых парах.}
\label{tab:ruler-interim-archive}
\footnotesize
\begin{tabular}{lrr}
\toprule
Показатель & 64 тыс. & 128 тыс. \\
\midrule
Сохранено пар из 390 & 368 & 225 \\
Оценка каждого варианта, \% & 98.91 & 100.00 \\
Полностью верно в каждом варианте & 364/368 & 225/225 \\
Макнемар по полностью верным ответам & 1 & 1 \\
Изменён текст: оффлоад & 28 (7.61\%) & 12 (5.33\%) \\
Изменён текст: другая разбивка & 32 (8.70\%) & 12 (5.33\%) \\
Только оффлоад / только разбивка & 12/16 & 7/7 \\
Макнемар по изменению текста & 0.5716 & 1 \\
Избыток изменений, п.п. & -1.09 & 0.00 \\
95\%-й интервал избытка, п.п. & [-5.83, +3.69] & [-5.72, +5.72] \\
\bottomrule
\end{tabular}
\end{table}

На 64 тысячах интервалы Уилсона для изменённого текста равны 5.32--10.78\% у оффлоада
и 6.23--12.02\% у другой разбивки; на 128 тысячах 3.08--9.09\% в обеих парах.
Макнемар по полностью верным ответам равен 1 отдельно для оффлоада с локальным и разбивки с локальным.
На 64 тысячах отсутствуют 22 примера \texttt{qa\_2}; на 128 тысячах сохранены только ранние задачи поиска иголок,
отсутствуют 165 примеров, включая все \texttt{vt/cwe/fwe/QA}.
Эти префиксы не устанавливают зарегистрированную эквивалентность.

На момент архивного среза исходные узлы были поставлены в очередь на перезапуск; завершение требовало восстановления замороженных спецификаций
и отдельной идентичности сборки и среды каждой глубины. Модель, набор данных, seed и исходный коммит совпадают,
но манифесты различаются по 64 хэшам нативных бинарных файлов.
Одних сохранённых строк недостаточно для автоматического продолжения.

\subsection{Грубая оценка начальной фазы}
\label{app:start-budget}
Для 100 тысяч KV q8\_0 требуется $100000\times34816=3481600000$ байт,
около 3.5~ГБ без рекуррентного состояния и обрамления.
Замороженная конструкция достигает около 100 тысяч за четыре запроса;
измеряемый вход начала равен 91\,214 токенам, из них 24\,378 переменных в первом запросе.
Грубая оценка только prefill использует измеренные 8122 токена/с при 16 тысячах
и порядок 3000--5000 токенов/с на 65--100 тысячах:
$24378/8122+(91214-24378)/5000=16.37$ секунды,
$24378/8122+(91214-24378)/3000=25.28$ секунды GPU.
Перенос скорости первой опоры на порцию и последующих скоростей на весь остаток есть допущения,
без измерения полного начального пути. Генерация, очереди, подготовка копии и передача в оценку GPU не входят.
Для M3 Max нормировка примерно на 200 токенов/с даёт $91214/200=456.07$ секунды,
7.60 минуты; чтение всех 100 тысяч дало бы $100000/200/60=8.33$ минуты.
Это порядок исходной оценки 6--9 минут, без нового замера.

Отношение $13169/233=56.52$ смешивает числитель full04 с активными датами full05;
такое отношение служит только грубой нормировкой, без наблюдённой частоты появления контекстов.
Согласованный full05 даёт $13175/233=56.55$ потока на активную дату, также без оценки новых задач.
Если принять около 56.5 новых контекстов в день как дополнительное предположение интенсивности,
то $56.5\times(16.37\ldots25.28)/60=15.4\ldots23.8$ минуты GPU на человека в день,
порядка 20 минут. Подагенты включены; это оценка при выбранной интенсивности,
без наблюдённого бюджета GPU и без сквозного прогноза зарегистрированного опыта.


\subsection{Трафик архитектуры и ёмкость начальной фазы}
\label{app:architecture-traffic}
Все объёмы и времена ниже являются арифметикой по измеренным размерам и скоростям;
передача четырёх начальных запросов отдельным сетевым опытом не измерена. ГБ и МБ десятичные.
Рассматриваем 100\,000 новых позиций и четыре передачи, без повторов, обрамления и подготовки.
Объём предполагает все четыре оффлоада, без решения политики маршрутизации.
Готовый общий префикс уменьшает переменную передачу; оценка полного объёма и оценка GPU ниже
с однократным вычислением общего префикса имеют разные границы учёта.
Qwen q8\_0 даёт $100000\times34816=3481600000$ байт KV;
состояние bf16 добавляет $4\times78446592=313786368$ байт,
всего $3795386368$ байт, около 3.8~ГБ. Число передач состояния является отдельным условием.

Glimmer в сохранённых слотах занимает 7.3--9.6~КиБ на токен;
отношение уже включает живое скользящее окно и данные слота.
Произведение $100000\times(7.3\ldots9.6)\times1024$ даёт 0.748--0.983~ГБ,
но удельный размер слота не является приростом сетевого объёма на новый токен.
Принятая геометрия отделяет 7072 байта глобального KV на позицию и 43\,450\,368 байт живого окна на обмен
(раздел~\ref{sec:models}). После заполнения окна четыре передачи дают
$100000\times7072+4\times43450368=881001472$ байта, около 0.9~ГБ.
Это сумма при передаче всего живого окна каждый раз, без измерения данного пути переноса;
повторное умножение удельного размера слота на число окон неверно.

\begin{table}[t]
\centering\footnotesize
\caption{Расчёт начальной фазы: 100 тысяч новых позиций и четыре обмена. Передаётся KV q8\_0;
Qwen добавляет состояние bf16, Glimmer всё живое окно. Времена делят объём на скорость, без RTT и накладных расходов.}
\label{tab:start-traffic}
\begin{tabular}{lrr}
\toprule
Величина & Qwen3.8-27B & Glimmer 30B \\
\midrule
Объём, ГБ & 3.795 & 0.881 \\
Время при 125 МБ/с, с & 30.4 & 7.0 \\
Время при 20 МБ/с, с & 189.8 & 44.1 \\
Время при 5 МБ/с, с & 759.1 & 176.2 \\
Десять начал в день, ГБ & 37.95 & 8.81 \\
\bottomrule
\end{tabular}
\end{table}

Для приближённой ёмкости GPU берём порции 24 / 35 / 20 / 12 тысяч токенов, всего 91 тысяча;
общий префикс 9624 считается один раз вне повторяющейся нагрузки.
Измеренные скорости Qwen vLLM для порции 10 тысяч равны 7144 / 4598 / 3130 токенов/с
на 16 / 65 / 131 тысяче; линейная интерполяция на представительных 50 / 80 / 90 тысячах
даёт 5377 / 4264 / 4042. Для прикидки округляем до 7.1 / 5.4 / 4.3 / 4.0 тысячи токенов/с:
\[
T_Q=24000/7100+35000/5400+20000/4300+12000/4000=17.51~\text{с}\approx18~\text{с}.
\]
Получаем $86400/18=4800$ начал в сутки и $0.7\times4800=3360\approx3400$ при 70\% занятости;
при десяти началах на человека это 336--480, округлённо 340--480 человек.
Предел относится только к GPU prefill при условии квалификации полного пути, без серверной генерации.
Для Glimmer используем собственные измерения llama.cpp: порция 10 тысяч обрабатывается при
4221 / 3712 / 3126 токенах/с на 16 / 50 / 100 тысячах.
Интерполяция на 80 / 95 тысячах даёт 3360.4 / 3184.6;
$24000/4221+35000/3712+20000/3360.4+12000/3184.6=24.83$ секунды.
Оценка по нижней скорости $91000/3000=30.33$ секунды задаёт порядок 25--30 секунд;
производительность vLLM для Glimmer не измерена.

Исходящий канал 1~Гбит/с при 125~МБ/с даёт
$86400\times125\cdot10^6=10.8\cdot10^{12}$ байт в сутки.
Qwen ограничен $10.8\cdot10^{12}/3795386368=2846$ началами, около 2800;
Glimmer при точном расчётном объёме $10.8\cdot10^{12}/881001472=12259$.
При округлённых 0.9~ГБ это 12\,000, при 1~ГБ 10\,800, около 11 тысяч.
Эти пределы канала и GPU не складываются: ограничивает меньший ресурс.
Поступление равномерно, тогда как реальные дни имеют пики; распределение очереди не измерено.
Контексты субагентов часто меньше медианного первого запроса 34 тысячи по характеристике владельца,
без измеренной доли таких контекстов. Общий префикс вычисляется один раз;
принудительное принятие десяти новых начал в день является сценарием интенсивности, без наблюдённой частоты.
